\section{CONCLUSION}\label{sec:conclusion}  
%%%%%%%%%%%%%%%%%%%%%%%%%%%%%%%%%%%%%%%
%% Add conclusions here
The interdisciplinary approach of this bachelor thesis to unit the small scaled investigation of the biochemistry with the greater view of the geography allowed a more holistic view towards the problem of hard-to-abate methane emissions from the agriculture sector. \\

The analysis of national methane emissions for the three countries of interest (\acrshort{deu}, \acrshort{nzl}, and the \acrshort{usa}) showed that the agriculture sector is a major potential source for reducing methane emissions. Therefore, as the dominant emission source within the agriculture sector, methane emissions from enteric fermentation processes in cattle are a valuable mitigation target. For \acrshort{deu} and \acrshort{nzl}, agriculture is the only sector left with potential for relevant methane reduction (see \cref{fig:sectors_emission}).

The general geodata analyses of the national methane emissions from \acrshort{deu}, \acrshort{nzl} and \acrshort{usa} proofed the assumption derived from global methane emission pattern, there is a need for a fitting mitigation strategy toward \ch{CH4} emission from enteric fermentation processes in cattle \parencite{appelhansUnterschatztesTreibhausgasMethan2025, saunoisGlobalMethaneBudget2025, IPPCC_AR6_Chapter7_GWP}. In addition, this single emission source contributes that much to the national emission (from 25~\% in the \acrshort{usa} 2022 to 45~\% in \acrshort{deu} and 60~\% in \acrshort{nzl}, \cref{fig:share_of_EnterFe}) it will have a relevant impact on mitigate climate change, if international addressed. Especially the calculation of the \acrshort{gwp}20 and \acrshort{gwp}100 in \cref{fig:saved_CH4} showed clearly, methane reduction could be really helpful short-term to stay at least below the 2\,°C aim of the \textcite{unfcccParisAgreement2015}. For this purpose, as already pointed out in \cref{sec:discuss_eval_immun} a shift of the rumen microbiome towards non-methanogenic microbes needs to be at least at 30~\%.
Wether the immunization of cattle can be helpful for this short-term mitigation, depends on the stage of biochemist development and investigation of a potent and specific vaccine against methanogenic \gls{archa}.

The biochemistry part of this bachelor thesis, could collect several basic information on the poorly understood genus \acrlong{mbrevi} as the dominant methanogenic \gls{archa} in the rumen of cattle \parencite[p.\,3]{baca-gonzalezAreVaccinesSolution2020}. Here, promising results were achieved to implement for the first time a genetic toolbox for \acrshort{bovi} and \acrshort{milli}. For this purpose, a transformation strategy was used based on interdomain conjugation with \acrshort{ecol}. If this is confirmed by whole-genome sequencing, it would allow a new approach for identifying a specific and effective antigen against \acrlong{mbrevi}, which would be a huge step towards immunization of cattle. Besides this, it could be shown that \acrshort{thau} and \acrshort{bovi} can be cultured and plated on media (DSM 1523 and 119) without rumenfluid and the \acrlong{mbrevi} are sensitive towards 20~µg/mL puromycin. In addition, first cellcounts per mL were determined (\cref{fig:cell_count_bovis_bsp}) allowing first calculation of plating efficiency. 

One problem has been identified and has to be resolved for future experiments, is that the solvent \acrshort{dmso} has a clear inhibitory effect on the growth of \acrshort{mbrevi} (\cref{sec:azaeight_result}) due to its oxidative effect \parencite{atcherEffectDMSOAqueous2013}. For chemicals that are poorly soluble in water, such as \acrshort{azaeight}, which is needed for the \gls{counterselection}, an alternative solvent must be used. 

As becomes obvious here, research on a possible vaccine for cattle against their methane-producing microbiome is still in the beginning, therefore some time is still needed until such a vaccine could become available on the market. This implies that the mitigation strategy will not help mitigate climate change within the next decade.
Therefore, the best approach to notably reduce methane over the next years is to use the reduction potential in the energy and waste sectors, as is clearly evident for the \acrshort{usa} (\cref{fig:sectors_emission}). As shown for \acrshort{deu} ()\cref{fig:sectors_emission}), these sectors are easier to address with today's technology \parencite{appelhansUnterschatztesTreibhausgasMethan2025, ipccClimateChange20222022}.

Nevertheless, this immunization strategy can be highly useful in the long term for achieving the goal of net-zero emissions within the agricultural sector, at least to some extent  \parencite{unfcccParisAgreement2015}. This will become even more relevant considering that methane emissions from natural sources, such as permafrost, are expected to increase in the future due to rising temperatures \parencite{schadelPermafrostWildfireCarbon2026}. However, it must be noted that even with a mitigation strategy reducing emissions per head of cattle, it will still be necessary to reduce the cattle population to a sustainable size in the future \parencite{edelenboschReducingSectoralHardtoabate2024}. This will be challenging given the growing demand for animal products \parencite{faoLivestockBalance2009, appelhansUnterschatztesTreibhausgasMethan2025}. Here, immunization or other strategies to reduce \ch{CH4} production in cattle could allow for a population size that is both socially and politically acceptable.

Future steps for the biochemical part are to further improve the \gls{trafoglo} protocol for \acrlong{mbrevi} and conduct initial genetic analyses. A highly interesting value to determine is the actual methane emission of \acrlong{mbrevi} via gas chromatography. For the geographical part, the geo-data analyses should be extended to other countries and it should also be investigated whether the mitigation strategy affects or interacts with other emission sectors, such as manure management or pasture carbon sinks.

At last, the basic research on methane production by methanogenic archaea conducted here, may be relevant not only for vaccine development, but also for other mitigation strategies, such as feed additives, and for our general understanding of methane production from natural sources.
